\section{Motivation}

% Ý đoạn: đặt bối cảnh lớn của bài toán AMR ở mức y tế công cộng.
% Vai trò: giải thích vì sao chủ đề này đáng nghiên cứu ngay từ đầu báo cáo.
Antimicrobial resistance (AMR) has become a major global public-health threat.\footnote{World Health Organization fact sheet on antibiotic resistance: \url{https://www.who.int/news-room/fact-sheets/detail/antibiotic-resistance}.} The World Health Organization notes that bacterial AMR was directly responsible for 1.27 million deaths in 2019 and contributed to 4.95 million deaths worldwide, which shows that resistance is not a narrow laboratory issue but a large-scale clinical and public-health problem. These burden estimates are also supported by a large global analysis published in \textit{The Lancet}~\citep{amrcollaborators2022global}.

% Ý đoạn: kéo vấn đề từ mức vĩ mô xuống mức thực hành lâm sàng và giám sát.
% Vai trò: cho thấy tại sao dự đoán phenotype kháng thuốc có ý nghĩa thực tế.
From a practical perspective, identifying whether a bacterial isolate is resistant or susceptible to a given antibiotic is essential for effective treatment, antibiotic stewardship, and surveillance. A wrong or delayed resistance assessment can lead to inappropriate therapy, prolonged infection, and wider spread of resistant strains. For this reason, resistance phenotype is not only a microbiological label; it is a clinically meaningful endpoint that directly affects how infections are managed~\citep{shen2024comparison}.

% Ý đoạn: nối nhu cầu thực tế với hướng tiếp cận của luận văn bằng dữ liệu genomic.
% Vai trò: chốt động lực trực tiếp của nghiên cứu và dẫn sang bài toán của report.
This practical need motivates the present study. Although antimicrobial susceptibility testing remains the reference standard, genomic data offers the possibility of inferring resistance from detected genes and mutations in a faster and more scalable way~\citep{feldgarden2021amrfinderplus,ren2022prediction}. The central motivation of this report is therefore to study whether curated AMR determinants can be transformed into a reliable antibiotic-specific prediction workflow that remains biologically interpretable while still achieving useful predictive performance.
